% Generated draft: requires mathematical and source review.
\documentclass[11pt]{article}
\usepackage[margin=25mm]{geometry}
\usepackage{fontspec}
\setmainfont{Latin Modern Roman}
\usepackage{amsmath,amssymb,mathtools,unicode-math}
\setmathfont{Latin Modern Math}
\usepackage{enumitem,needspace,xurl,hyperref}
\setlength{\parindent}{0pt}
\setlength{\parskip}{5pt}
\setlength{\emergencystretch}{3em}
\newcommand{\N}{\mathbb{N}}
\newcommand{\Z}{\mathbb{Z}}
\newcommand{\Q}{\mathbb{Q}}
\newcommand{\R}{\mathbb{R}}
\newcommand{\C}{\mathbb{C}}
\newcommand{\F}{\mathbb{F}}
\newcommand{\A}{\mathbb{A}}
\newcommand{\PP}{\mathbb{P}}
\newcommand{\E}{\mathbb{E}}
\newcommand{\eps}{\varepsilon}
\newcommand{\dd}{\,\mathrm{d}}
\begin{document}
\section*{UnsolvedMath ID: 30004340}
\textbf{Source problem code:} catalogue entry
\subsection{Definitions and mathematical statement}
Let \(N\in\N\) denote the matrix size, and let \(W=W(N)\) denote the bandwidth parameter of a sequence of random band matrices \(H_N=(h_{ij}^{(N)})_{1\leq i,j\leq N}\). The term random band matrix indicates that the matrix interactions are restricted to, or concentrated near, a band around the diagonal; the supplied record does not specify which precise convention is intended.

\textbf{Requested classification.} Determine the universality classes of these matrices as \(N\to\infty\) and the bandwidth \(W(N)\) varies, including the behavior across the localization--delocalization transition.

\textbf{Unresolved specifications.} This is a research question rather than a uniquely specified mathematical conjecture. The record does not fix the matrix symmetry class, the index geometry and boundary conditions, the bandwidth convention or admissible bandwidth sequences, the joint distribution and dependence of the entries, their centering and variance normalization, or their moment assumptions. It also does not identify the observables, spectral region, rescaling, mode of convergence, or equivalence criterion that define the requested universality classes. No operational definition of localization or delocalization, transition scale, or critical scaling regime is supplied. These choices must be specified before a precise limiting assertion can be formulated; none is selected here.
\subsection{Short English statement}
The question asks how the universal limiting behavior of random band matrices depends on bandwidth as matrix size grows, including through the localization--delocalization transition. The supplied record leaves the ensemble, observable, scaling, and meaning of the transition unspecified, so it does not determine a single precise conjecture.
\subsection{Sources}
\begin{itemize}
\item Supplied local UnsolvedMath registry record: permanent ID \(30004340\), source problem code \texttt{catalogue entry}, title \emph{Universality Across the Random-Band-Matrix Localization Transition}; collection: Oberwolfach Reports Open Problems. The supplied proposed year is 2019; no individual proposer is identified.
\item Supplied report attribution: \emph{Random Matrices (2020)}, Oberwolfach Reports, DOI 10.4171/owr/2019/56, pp.\ 3459--3527. \url{https://doi.org/10.4171/owr/2019/56}. The background describes this collection as drawn from workshop reports and their open-problem discussions; no problem-specific page is supplied.
\item Supplied literature source, labeled ``OWR report on random matrices'': \url{https://oa.tib.eu/renate/bitstreams/bd70698e-b340-45c1-9989-9cb604446b31/download}.
\item Supplied literature source, labeled ``Bourgade--Erdos--Yau--Yin, universality for random band matrices'': \url{https://arxiv.org/abs/1602.02312}.
\item Supplied literature source, labeled ``recent delocalization results for random band matrices'': \url{https://arxiv.org/abs/2501.01718}.
\item Supplied literature assessment dated \texttt{2026-08-22}: the record labels the problem partially solved and reports that universality is established in several bandwidth regimes, including mean-field regimes and widths above critical thresholds, while the full varying-width universality classification and transition remain open. This is supplied dated curation derived from the OWR material, not an independently verified literature assessment. The source itself requests independent verification of time-sensitive status and cited claims; it supplies no numerical thresholds or applicable hypotheses here.
\item The supplied statement-status label is \texttt{corrected\_verified}. Its accompanying explanation says that mathematical content was preserved while notation, typography, omitted definitions, or duplicate clauses were repaired for self-contained publication. This is registry metadata, not independent mathematical verification of this draft.
\end{itemize}
\end{document}
